\chapter{Introduction}
\label{ch:introduction}

This chapter sets out the problem, the gap in the literature, and the
research questions addressed by the thesis. It demonstrates the elements
the template provides: citations, cross-references, an equation, a table,
a figure, and acronyms.

\section{Motivation}

An \gls{abm} represents individual agents whose local interactions produce
system-level behaviour \autocite{bonabeau2002}. Such models are used where
analytical treatment is intractable, for example in evacuation studies
\autocite{helbing2000}.

\section{Research questions}

The thesis answers three questions:

\begin{enumerate}
  \item \rqref{1}: \rqone
  \item \rqref{2}: \rqtwo
  \item \rqref{3}: \rqthree
\end{enumerate}

\section{A quantified statement}

Let $R$ denote the evacuation ratio. We estimate it as
\begin{equation}
  R_{\mathrm{evac}} = \sum_{i=1}^{N} \frac{w_i\,\rho_i}{\tau_i + \varepsilon},
  \label{eq:revac}
\end{equation}
where $w_i$ is the weight of agent $i$, $\rho_i$ its density contribution,
$\tau_i$ its exit time, and $\varepsilon$ a small positive constant that
keeps the summand finite. \Cref{eq:revac} is referenced automatically.

\section{Illustrative results}

\Cref{tab:params} lists the parameters used throughout, and
\cref{fig:decay} shows the corresponding decay curve.

\begin{table}[htbp]
  \centering
  \caption{Model parameters.}
  \label{tab:params}
  \begin{tabular}{lrr}
    \toprule
    Parameter & Symbol & Value \\
    \midrule
    Number of agents & $N$ & 250 \\
    Exit width (m)   & $d$ & 1.2 \\
    Time step (s)    & $\Delta t$ & 0.25 \\
    \bottomrule
  \end{tabular}
\end{table}

\begin{figure}[htbp]
  \centering
  \begin{tikzpicture}[x=1.1cm,y=2.2cm]
    \draw[NottinghamBlue, thick, ->] (0,0) -- (6.2,0) node[right] {$t$};
    \draw[NottinghamBlue, thick, ->] (0,0) -- (0,1.35) node[above] {$R(t)$};
    \draw[MalaysiaSkyBlue, ultra thick, domain=0:6, samples=80]
      plot (\x, {exp(-0.45*\x)});
    \draw[UoNRed, dashed, thick] (0,1) -- (6,1) node[right, black] {initial};
  \end{tikzpicture}
  \caption{Evacuation ratio decaying over time.}
  \label{fig:decay}
\end{figure}

\section{Outline}

\Cref{ch:literature} reviews prior work, \cref{ch:methodology} describes
the method, \cref{ch:results} reports results, and \cref{ch:conclusion}
concludes.
